\subsection{Operaciones con matrices}

\textit{Igualdad}, dos matrices son iguales si esta tiene la misma dimension y elementos identicos en las ubicaciónes correspondientes. Es decir  \(A = [a_{ij}] = B = [b_{ij}]\) si \([a_{ij}] = [b_{ij}]\) para todos los valores de \(i\) y \(j\). Ejemplo 
\[
\begin{bmatrix}
    4 & 3 \\
    2 & 0
\end{bmatrix}
= 
\begin{bmatrix}
    4 & 3 \\
    2 & 0
\end{bmatrix}
\neq
\begin{bmatrix}
    2 & 0 \\
    4 & 3 
\end{bmatrix}
\]
Por lo tanto, una matriz 
\[
\begin{bmatrix}
    x \\
    y
\end{bmatrix}
= \begin{bmatrix}
    7 \\
    4
\end{bmatrix}
\]
si \(x = 7\) y \(y = 4\)

\subsubsection{Suma y resta de matrices}

Las matrices se pueden sumar si tienen la misma dimensión. Por lo tanto la suma de \(A = [a_{ij}]\) y \(B = [b_{ij}]\) es la suma de cada \textit{par de elementos correspondientes}
\begin{ejemplo}
    \[
    \begin{bmatrix}
        4 & 9 \\
        2 & 1
    \end{bmatrix}
    + 
    \begin{bmatrix}
        2 & 0 \\
        0 & 7 
    \end{bmatrix}
    = 
    \begin{bmatrix}
        4+2 & 9+0 \\
        2+0 & 1+7
    \end{bmatrix}
    = 
    \begin{bmatrix}
        6 & 9 \\
        2 & 8
    \end{bmatrix}
    \]
\end{ejemplo}
Por lo tanto una suma general se puede expresar como
\[
\begin{bmatrix}
    a_{11} & a_{12} \\
    a_{21} & a_{22} \\
    a_{31} & a_{32}
\end{bmatrix}
+
\begin{bmatrix}
    b_{11} & b_{12} \\
    b_{21} & b_{22} \\
    b_{31} & b_{32}    
\end{bmatrix}
= 
\begin{bmatrix}
    a_{11}+b_{11} & a_{12}+b_{12} \\
    a_{21}+b_{21} & a_{22}+b_{22} \\
    a_{31}+b_{31} & a_{32}+b_{32}
\end{bmatrix}
\]
Es decir que \([a_{ij}]+[b_{ij}]=[c_{ij}]\) donde \(c_{ij} = a_{ij} + b_{ij}\).

La matiz resultante debe tener la misma dimension que las matrices que la componen.

En el caso de la resta, es una operación suma, donde \(A + (-B)\) es la resta de matrices. La operacion se puede hacer si tienen la misma dimensión.

Entonces en el caso de la resta \([a_{ij}]-[b_{ij}]=[d_{ij}]\) donde \(d_{ij} = a_{ij} - b_{ij}\).

\begin{ejemplo}
    \[
    \begin{bmatrix}
        19 & 3 \\
        2 & 0 
    \end{bmatrix}
    - 
    \begin{bmatrix}
        6 & 8 \\
        1 & 3 
    \end{bmatrix}
    = 
    \begin{bmatrix}
        19-6 & 3-8 \\
        2-1 & 0-3 
    \end{bmatrix}
    = 
    \begin{bmatrix}
        13 & -5 \\
        1 & -3
    \end{bmatrix}
    \]
\end{ejemplo}

\subsubsection{Multiplicación escalar}

Un escalar es una magnitud que se define únicamente por un valor numérico y sus unidades, sin dirección ni sentido. Cuando multiplicamos una matriz por un número, estamos multiplicando cada elemento de esa matriz por ese número o escalar.
\begin{ejemplo}
    \[
    7\begin{bmatrix}
        3 & -1 \\
        0 & 5
    \end{bmatrix}
    = 
    \begin{bmatrix}
        7\times3 & 7\times-1 \\
        7\times0 & 7\times5
    \end{bmatrix}
    =
    \begin{bmatrix}
        21 & -7 \\
        0 & 35
    \end{bmatrix}
    \]
\end{ejemplo}
Un escalar tambien puede ser una fracción.

Por lo tanto la multiplicación por un escalar se puede representar de forma general como 
\[
n\begin{bmatrix}
    a_{11} & a_{12} \\
    a_{21} & a_{22}
\end{bmatrix}
= 
\begin{bmatrix}
    n\times a_{11} & n\times a_{12} \\
    n\times a_{21} & n\times a_{22}
\end{bmatrix}
\]
El escalar entonces incrementa o disminuye la matriz por cierto múltiplo.

Un escalar tambien puede ser un número negativo 
\begin{ejemplo}
    \[
    -1\begin{bmatrix}
        a_{11} & a_{12} & d_1\\
        a_{21} & a_{22} & d_2
    \end{bmatrix}
    = 
    \begin{bmatrix}
        -1\times a_{11} & -1\times a_{12} & -1\times d_1 \\
        -1\times a_{21} & -1\times a_{22} & -1\times d_2
    \end{bmatrix}
    = 
    \begin{bmatrix}
        -a_{11} & -a_{12} & -d_1\\
        -a_{21} & -a_{22} & -d_2
    \end{bmatrix}
    \]
\end{ejemplo}
Si la matriz representa los coeficientes y los términos constantes de ecuaciones simultáneas 
$$
\begin{aligned}
    & a_{11}x_1 + a_{12}x_2 = d_1 \\
    & a_{21}x_1 + a_{22}x_2 = d_2
\end{aligned}
$$
entonces la multiplicación por el escalar \(-1\) equivale a muiltiplicar ambos lados de ambas ecuaciones por \(-1\), y de este modo cambia el signo de todo el sistema.

\subsubsection{Multiplicación de matrices}

Para multiplicar matrices se debe satisfacer un requerimiento de dimension. Suponga una matriz \(A\) y una matriz \(B\), donde 
\[
A = 
\begin{bmatrix}
    a_{11} & a_{12}    
\end{bmatrix}
\quad
\text{y} 
\quad
B = 
\begin{bmatrix}
    b_{11} & b_{12} & b_{13} \\
    b_{21} & b_{22} & b_{23}
\end{bmatrix}
\]

La dimensión de \(A\) es de \(1 \times \color{red}2\) y la dimensión de \(B\) es de \({\color{red}2} \times 3\). El valor \(2\) en color rojo no es solo por gusto del editor, este valor es el importante. Una matriz se puede multiplicar cuando \(A\) contiene \(n\) columnas que son iguales a \(B\) \(m\) filas. Es decir una matriz se puede multiplicar cuando la dimensión de la \textit{fila} de \(A\) es igual a la dimensión de la\textit{columna} de \(B\).

Por lo tanto el producto \(AB\) se puede multiplicar, pero si quisieramos multiplicar \(BA\) note la diferencia
\[
B = 
\begin{bmatrix}
    b_{11} & b_{12} & b_{13} \\
    b_{21} & b_{22} & b_{23}
\end{bmatrix}
\quad
\text{y} 
\quad
A = 
\begin{bmatrix}
    a_{11} & a_{12}    
\end{bmatrix}
\]
Las dimensiones siguen siendo iguales, pero ya no satisfacemos la igualdad de el numero de filas y columnas, por lo tanto no se puede multiplicar, ya que \(B = 2 \times {\color{red}3}\) y \(A = {\color{red}1} \times 2 \). A esta condición se le llama condición de \textit{conformabilidad}.

\begin{nota}
    Para que \(B\) se pueda multiplicar por \(A\), es decir \(BA\), \(B\) deberia tener 1 fila, o \(A\) tener \(3\) columnas.
\end{nota}

El producto de \(AB\) esta definido y se espera que sea de la dimensión \(1 \times 3\), por lo tanto 
\[
AB = C =
\begin{bmatrix}
    c_{11} & c_{12} & c_{13} 
\end{bmatrix}
\]
Para calcular \(c_{11}\), estamos haciendo el siguiente proceso \(a_{11}\times b_{11} + a_{12} \times b_{21}\). En el caso de \(c_{12}\), \(a_{11}\times b_{12} + a_{12} \times b_{22}\) y por ultimo para \(c_{13}\) \(a_{11}\times b_{13} + a_{12} \times b_{23}\).

Note que para calcular \(c_{12}\) se multiplica \(a_{1{\color{red}1}}\) por \(b_{{\color{red}1}2}\) y se le suma la multiplicación de \(a_{1{\color{red}2}}\) por \(b_{{\color{red}2}2}\). Para \(c_{11}\) se ilustra en el grafico \ref{Multiplicacion_de_matrices}

\begin{figure}[ht]
\centering
\begin{tikzpicture}[
    cell/.style={circle, draw, minimum size=8mm, inner sep=0pt},
    light/.style={fill=gray!20},
    dark/.style={fill=gray!50},
    bracket/.style={line width=0.8pt}
]
    %%% MATRIZ FILA (A) %%%
    \node[cell, light] (a11) at (0,0) {$a_{11}$};
    \node[cell, dark] (a12) [right=0.5cm of a11] {$a_{12}$};
    
    \draw[bracket] ([xshift=-2mm, yshift=4mm]a11.west) -- ([xshift=-4mm, yshift=4mm]a11.west) -- ([xshift=-4mm, yshift=-4mm]a11.west) -- ([xshift=-2mm, yshift=-4mm]a11.west);
    \draw[bracket] ([xshift=2mm, yshift=4mm]a12.east) -- ([xshift=4mm, yshift=4mm]a12.east) -- ([xshift=4mm, yshift=-4mm]a12.east) -- ([xshift=2mm, yshift=-4mm]a12.east);

    %%% MATRIZ COLUMNA (B) %%%
    \node[cell, light] (b11) at (5,0.5) {$b_{11}$};
    \node[cell, dark] (b21) [below=0.5cm of b11] {$b_{21}$};
    
    \draw[thick] (b11.east) ++(0.4,0) -- ++(0.4,0) (b11.east) ++(1.0,0) -- ++(0.4,0);
    \draw[thick] (b21.east) ++(0.4,0) -- ++(0.4,0) (b21.east) ++(1.0,0) -- ++(0.4,0);

    % Corchete IZQUIERDO de B
    \draw[bracket] ([xshift=-4mm, yshift=4mm]b11.north west) -- ([xshift=-6mm, yshift=4mm]b11.north west) -- ([xshift=-6mm, yshift=-4mm]b21.south west) -- ([xshift=-4mm, yshift=-4mm]b21.south west);
    
    % Corchete DERECHO de B (MODIFICADO: xshift de 12/14 a 18/20)
    \draw[bracket] ([xshift=18mm, yshift=4mm]b11.north east) -- ([xshift=20mm, yshift=4mm]b11.north east) -- ([xshift=20mm, yshift=-4mm]b21.south east) -- ([xshift=18mm, yshift=-4mm]b21.south east);

    %%% CONEXIONES (ARCOS) %%%
    \draw[thick] (a11.north) .. controls +(up:1.5cm) and +(up:1.5cm) .. node[above, font=\small] {Primer par} (b11.north);
    
    % Segundo par (MODIFICADO: up:0.8cm a up:1.2cm para que suba)
    \draw[thick] (a12.north) .. controls +(up:1.2cm) and +(left:0.8cm) .. node[above, font=\small, xshift=0.3cm, yshift=0.3cm] {Segundo par} (b21.west);

    \node[anchor=west, font=\bfseries] at (-1.5, 1.5) {Para $c_{11}$:};
\end{tikzpicture}
\caption{Multiplicación de matrices, proceso}
\label{Multiplicacion_de_matrices}
\end{figure}

Es el requerimiento de formar parejas en este proceso el que requiere la comparación de la dimensión columna de la matris primaria y la dimensión de fila de la matriz secundaria.

El procedimiento de la multiplicación se puede describir por medio del concepto del producto interior de dos vectores. Dad un vector \(u\) y un vector \(v\) con \(n\) componentes cada uno, estos ordenados como dos filas y dos columnas o una fila y una columna su producto interior, va a ser escrito como \(u \centerdot v\) y se define como 
\[
u \centerdot v = u_1v_1 + u_2v_2 + \dots + u_nv_n
\] 
El producto interior de dos vectores es un escalar porque es una operación de agregación. Transforma dos listas de información en un único valor de síntesis que representa la intensidad de su interacción. Vease el ejemplo \ref{producto_interior_de_vectores}
\begin{ejemplo}
    \label{producto_interior_de_vectores}
    Después de un viaje de compras se ordenan las cantidades compradas de \(n\) bienes como un vector renglón 
    \(Q = \begin{bmatrix}
    Q_1 & Q_2 & \dots & Q_n
    \end{bmatrix}\)
    , y la lista de esos bienes en un vector precio \(
    P = \begin{bmatrix}
        P_1 & p_2 & \dots & P_n
    \end{bmatrix}
    \)
    entonces el producto interior de los vectores es 
    \[
    Q \centerdot P = Q_1P_1 + Q_2P_2 + \dots + Q_nP_n = \text{costo de compra total}
    \]
    Entonces si en ves de usar dos vectores lo hacemos con dos matrices es como pasar muchos carritos de compra por muchas cajas registradoras al mismo tiempo.Donde cada fila de la primera matriz \(A\) es un carrito de compras con una lista de cantidades \(Q\) y cada columna de la segunda matriz \(B\) es una lista de precios \(P\) de una tienda distinta. Por lo tanto el producto \(AB\) nos dara cuánto dinero \(c\) genero o gasto el carrito al pasar por la caja registradora.
\end{ejemplo}
\begin{ejemplo}
    Dado 
    \[
    A = \begin{bmatrix}
        1 & 3 \\
        2 & 8 \\ 
        4 & 0
    \end{bmatrix}
    \quad
    \text{y}
    \quad
    B = \begin{bmatrix}
        5 \\
        9
    \end{bmatrix}    
    \]
    determine \(AB\). El producto \(AB\) está definido porque \(A\) tiene dos columnas y \(B\) dos filas. Su matriz va a ser de \(3 \times 1\), por lo tanto
    \[
    AB = \begin{bmatrix}
        1 \times 5 + 3 \times 9 \\
        2 \times 5 + 8 \times 9 \\
        4 \times 5 + 0 \times 9 
    \end{bmatrix}
    =
    \begin{bmatrix}
        32 \\
        82 \\
        20
    \end{bmatrix}
    \]
\end{ejemplo}

\begin{ejemplo}
    \label{ejemplo_multipliación_matices_9}
    Dado 
    \[
    A = \begin{bmatrix}
        3 & -1 & 2 \\
        1 & 0 & 3 \\
        4 & 0 & 2 
    \end{bmatrix}
    \quad
    \text{y}
    \quad
    B = \begin{bmatrix}
        0 & -\frac{1}{5} & \frac{3}{10} \\
        -1 & \frac{1}{5} & \frac{7}{10} \\
        0 & \frac{2}{5} & -\frac{1}{10}
    \end{bmatrix}
    \]
    Obtenga \(AB\)
    $$
    \begin{aligned}
        & c_{11} = 3 \times 0 + (-1 \times -1) + 2 \times 0 &&= 1 \\
        & c_{12} = 1\times 0 + 0 \times -1 + 3 \times 0 &&= 0 \\
        & c_{13} = 4 \times 0 + 0 \times -1 + 2 \times 0 &&= 0 \\ 
        & c_{21} = 3 \times -\frac{1}{5} + (-1 \times \frac{1}{5}) + 2 \times \frac{2}{5} &&= 0 \\
        & c_{22} = 1\times -\frac{1}{5} + 0 \times \frac{1}{5} + 3 \times \frac{2}{5} && = 1 \\
        & c_{23} = 4 \times -\frac{1}{5} + 0 \times - \frac{1}{5} + 2 \times \frac{2}{5} && = 0 \\ 
        & c_{31} = 3 \times  \frac{3}{10} + (-1 \times \frac{7}{10}) + 2 \times -\frac{1}{10} && = 0 \\
        & c_{32} = 1\times \frac{3}{10} + 0 \times \frac{7}{10} + 3 \times -\frac{1}{10} && = 0 \\
        & c_{33} = 4 \times \frac{3}{10} + 0 \times \frac{7}{10} + 2 \times -\frac{1}{10} && = 1 \\
    \end{aligned}
    $$
    Por lo tanto en forma de matriz el resultado de \(AB\) seria: 
    \[
    AB = \begin{bmatrix}
        1 & 0 & 0 \\
        0 & 1 & 0 \\
        0 & 0 & 1
    \end{bmatrix}
    \]

    Esta matriz ademas tiene algo curioso, a esta matriz se le llama matriz identidad, ya que su diagonal principal tiene valores pero su traingulo inferior y superior vale \(0\).
\end{ejemplo}
\begin{ejemplo}
    Suponga una matriz \(A\) y un vector \(x\):
    \[
    A = \begin{bmatrix}
        6 & 3 & 1 \\
        1 & 4 & -2 \\
        4 & -1 & 5 
    \end{bmatrix}
    \times
    x = \begin{bmatrix}
        x_1 \\
        x_2 \\ 
        x_3
    \end{bmatrix}
    Ax = \begin{bmatrix}
        6x_1 + 3x_2 + x_3 \\
        x_1 + 4x_2 +- 2x_3 \\
        4x_1 - x_2 + 5x_3
    \end{bmatrix}
    \]
    El producto es de \(3 \times 1\). 

    El resultado de esta matriz nos daria un vector columna \(d\) con los valores de solución
\end{ejemplo}
\begin{ejemplo}
    En el modelo de ingreso nacional simple de dos variables endogenas \(Y\) y \(C\), donde 
    $$
    \begin{aligned}
        & Y = C + I_0 + G_0 \\
        & C = a + bY
    \end{aligned}
    $$
    Si despejamos los terminos independientes \((I_0, G_0, a)\)
    $$
    \begin{aligned}
        & Y - C = I_0 + C_0 \\ 
        & -bY + C = a
    \end{aligned}
    $$
    Si lo pasamos a matriz, los coeficientes se encontrarian en \(A\), las variables en \(x\) y las constantes en \(d\), que quedaria 
    \[
    A = \begin{bmatrix}
        1 & -1 \\
        -b & 1 
    \end{bmatrix}
    x = \begin{bmatrix}
        Y \\
        C
    \end{bmatrix}
    d = \begin{bmatrix}
        I_0+G_0 \\
        a
    \end{bmatrix}
    \]
    Aplicando la regla de multiplicación de matrices se tiene que 
    \[
    Ax = \begin{bmatrix}
        1 \times Y + (-1) \times C \\ 
        -b \times Y + 1 \times C
    \end{bmatrix} 
    = 
    \begin{bmatrix}
        I_0+G_0 \\
        a
    \end{bmatrix}
    \]
    Por lo tanto la ecuación matricial \(Ax = d\) produciria 
    \[
    Ax = \begin{bmatrix}
        Y - C \\ 
        -bY + C
    \end{bmatrix} 
    = 
    \begin{bmatrix}
        I_0+G_0 \\
        a
    \end{bmatrix}
    \]
    Como la igualdad de matrices indica la igualdad entre elementos, es claro que la ecuación \(Ax = d\) representa con presición el sistema de ecuaciones original. 
\end{ejemplo}




\subsubsection{División}

La división de matrices no existe como tal. En el caso de la division de dos números \((a, b)\) la division de \(a/b\) es lo mismo que decir \(ab^{-1}\) donde \(b^{-1}\) representa el \textit{inverso} o \textit{reciproco}. Como \(ab^{-1}\) y \(b^{-1}a\) son iguales, la expresion del cociente \(a/b\) se puede representar tanto como \(ab^{-1}\) y \(b^{-1}a\).

Pero esto en las matrices es diferente, para aplicar el concepto de inverso en matriz, una matriz \(B^{-1}\) es la matriz inversa de \(B\), pero la condición de conformabilidad deduce que \(AB^{-1}\) puede estar definido, pero puede que \(B^{-1}A\) no estar definido, y en el caso que estuvieran definidos pueden no representar el mismo producto.

Por lo tanto \(AB^{-1}\) o \(B^{-1}A\) se diferencian, requiriendo que se especificque, siempre que exista el inverso de \(B\) se puede dividir (y que cumpla con el requisito de conformabilidad.)


\subsubsection{Notación \(\Sigma\)}

Los simbolos subíndice \((x_n)\) no son solo para designar el lugar del parámetro y variables, son también una abreviatura flexible que dentoa suma de términos, como los que surgen durante el proceso de multiplicación de matrices.

La abreviatura de suma utiliza \(\Sigma\) (sigma) como representacion de la suma. Para expresar la suma de \(x_1, x_2\) y \(x_3\) se puede escribir 
\[
x_1 + x_2 + x_3 = \sum_{j=1}^{3} x_j
\]
Esto se lee como ``La suma de \(x_j\) cuando \(j\) varía de \(1\) a \(3\)''. \(j\) es conocido como el \textit{índice de sumatoria}, toma sólo valores enteros. La expresión \(x_j\) representa al \textit{sumado}. Ademas de \(j\), los índices de la sumatoria también se denotan como \(i\) o \(k\) como 
$$
\begin{aligned}
    & \sum_{i=3}^{7} x_i =x_3 + x_4 + x_5 + x_6 +x_7 \\
    & \sum_{k=0}^{n} x_i =x_0 + x_1 + \dots + x_n 
\end{aligned}
$$
La aplicación de la notación suma se puede extender a los casos en los que el término \(x\) lleva como prefijo un coeficiente o en los que cada término de la suma se eleva a una potencia entera. Por ejemplo 
$$
\begin{aligned}
    & \sum_{j=1}^{3} ax_j = ax_1 + ax_2 + ax_3 = a(x_1 + x_2 + x_3) = a \sum_{j=1}^{3} x_j \\
    & \sum_{j=1}^{3} a_jx_j = a_1x_1 + a_2x_2 + a_3x_3 \\
    & \sum_{i=0}^{n} a_ix^i = a_0x^0 + a_1x^1 + a_2x^2 + \dots + a_nx^n = a_0 + a_1x + a_2x^2 + \dots + a_nx^n 
\end{aligned}
$$
En el último caso se muestra \(\sum_{i=0}^{n} a_ix^i\) que se puede usar de forma abreviada para representar una función polinomial general.

Siempre que no deje ambigüedad se puede omitir el alcance de la suma, el simbolo \(\sum\) se puede usar solo, o sólo con la letra índice de abajo.

Cada elemento de la matriz producto \(C = AB\) se define como una suma de término que se puede reescribir como 
$$
\begin{aligned}
    & c_{11} = a_{11}b_{11} + a_{12}b_{21} = \sum_{k=1}^{2} a_{1k}b_{k1} \\
    & c_{12} = a_{11}b_{12} + a_{12}b_{22} = \sum_{k=1}^{2} a_{1k}b_{k2} \\
    & c_{13} = a_{11}b_{13} + a_{12}b_{23} = \sum_{k=1}^{2} a_{1k}b_{k3}
\end{aligned}
$$

En la multiplicación de matriz \(A = [a_{ik}]\) de \(m \times n\) y de una matriz \(B=[b_{kj}]\) de \(n \times p\), se puede escribir los elementos de la matriz producto \(AB = C = [c_{ij}]\) de \(m \times p\) como
\[
c_{11} - \sum_{k=1}^{n} a_{1k}b_{k1} \qquad  c_{12} - \sum_{k=1}^{n} a_{1k}b_{k2}
\]
o de modo general 
\[
c_{ij} = \sum_{k=1}^{n} a_{ik}b_{kj} \qquad \begin{matrix}
    i = 1,2, \dots, m \\
    j = 1,2, \dots, p
\end{matrix}
\]

Esta última ecuación representa otra forma de expresar la multiplicación de matrices definidas.

\begin{ejercicio}
    \begin{enumerate}
        \item Dado \(A, B\) y \(C\) obtenga: 
        \((a)\) \(A+B\) \quad \((b)\) \(C - A\) \quad \((c)\) \(3A\) \quad \((d)\) \(4B = 2C\)
        Donde 
        \[
        A = \begin{bmatrix}
            7 & -1 \\
            6 & 9 
        \end{bmatrix}
        , B= \begin{bmatrix}
            0 & 4 \\
            3 & -2 
        \end{bmatrix} 
        \text{ y }
        C = \begin{bmatrix}
            8 & 3 \\
            6 & 1
        \end{bmatrix}
        \]
        \item Dado  
        \[
        A = \begin{bmatrix}
            2 & 8 \\ 
            3 & 0 \\
            5 & 1 
        \end{bmatrix}
        , B = \begin{bmatrix}
            2 & 0 \\
            3 & 8 
        \end{bmatrix}
        \text{ y }
        C = \begin{bmatrix}
            7 & 2 \\
            6 & 3
        \end{bmatrix}
        \] 
        \((a)\) ¿Está definido \(AB\)? Calcule \(AB\) ¿Puede calcular \(BA\)? ¿Por qué?
        \((b)\) ¿Está definido \(BC\)? Calcule \(BC\) ¿Está definido \(CB\)? En caso afirmativo, calcule \(CB\) ¿Es cierto que \(BC =CB\)?
        \item Con base en las matrices dadas en el ejemplo \ref{ejemplo_multipliación_matices_9}, las cuales eran:
            \[
            A = \begin{bmatrix}
                3 & -1 & 2 \\
                1 & 0 & 3 \\
                4 & 0 & 2 
            \end{bmatrix}
            \quad
            \text{y}
            \quad
            B = \begin{bmatrix}
                0 & -\frac{1}{5} & \frac{3}{10} \\
                -1 & \frac{1}{5} & \frac{7}{10} \\
                0 & \frac{2}{5} & -\frac{1}{10}
            \end{bmatrix}
            \]
        ¿Está definido el producto \(BA\)? Si es así,calcule el producto. ¿En este caso se tiene \(AB = BA\)? 
        \item Obtenga las matrices producto de los siguentes casos 
        $$
        \begin{aligned}
        & (a) \begin{bmatrix}
            0 & 2 & 0 \\
            3 & 0 & 4 \\
            2 & 3 & 0
        \end{bmatrix}
        \begin{bmatrix}
            8 & 0 \\
            0 & 1 \\
            3 & 5
        \end{bmatrix}
        \qquad
        && (b) \begin{bmatrix}
            6 & 5 & -1 \\
            1 & 0 & 4
        \end{bmatrix}
        \begin{bmatrix}
            4 & -1 \\
            5 & 2 \\
            0 & 1
        \end{bmatrix} \\
        & (c) \begin{bmatrix}
            3 & 5 & 0 \\
            4 & 2 & -7
        \end{bmatrix}
        \begin{bmatrix}
            x \\
            y \\
            z
        \end{bmatrix}
        \qquad
        && (d) \begin{bmatrix}
            a & b & c
        \end{bmatrix}
        \begin{bmatrix}
            7 & 0 \\
            0 & 2 \\
            1 & 4
        \end{bmatrix}
        \end{aligned}
        $$ 
        \item En el ejemplo \ref{producto_interior_de_vectores} si se dispone de las cantidades y precios como vectores columna en lugar de vectores filas ¿Está definido \(Q \centerdot P\)?
        \item Desarrolle las siguientes expresiones de suma  
        $$ 
        \begin{aligned}
            & (a) \sum_{i=2}^{5} x_i  && (b) \sum_{i=5}^{8} a_ix_i \\
            & (c) \sum_{i=1}^{4} bx_i && (d) \sum_{i=1}^{n} a_ix^{i-1} \\
            & (e) \sum_{i=0}^{3} (x+i)^2
        \end{aligned}
        $$
        \item Reescriba lo siguiente en notación de \(\sum\) 
        $$ 
        \begin{aligned}
            & (a) x_1(x_1 - 1) + 2x_2(x_2 - 1) + 3x_3(x_3 - 1) \\
            & (b) a_2(x_3+2) + a_3(x_4+3)+a_4(x_5+4) \\
            & (c) \frac{1}{x}+ \frac{1}{x^2} + \dots + \frac{1}{x^n} \quad (x \neq 0) \\
            & (d) 1 + \frac{1}{x} + \frac{1}{x^2} + \dots + \frac{1}{x^n} \quad (x \neq 0) 
        \end{aligned}
        $$
        \item Muestre que las siguientes expresiones son ciertas
        $$ 
        \begin{aligned}
            & (a) \left(\sum_{i=0}^{n} x_i\right) + x_{n+1} = \sum_{i=0}^{n+1} x_i \\
            & (b) \sum_{j=1}^{n} ab_j y_j = a \sum_{j=1}^{n} b_j y_j \\
            & (c) \sum_{j=1}^{n} (x_j+y_j) = \sum_{j=1}^{n} x_j + \sum_{j=1}^{n} y_j
        \end{aligned}
        $$
    \end{enumerate}
\end{ejercicio}